\chapter{Fraud, Risk, and Trust \& Safety (Week 10)}

\section{Learning objectives}
\begin{itemize}
  \item Identify major categories of fraud and abuse in e-commerce, including payment fraud, account takeover, promotion abuse, and refund abuse.
  \item Explain why fraud detection is a cost-sensitive, adversarial, and highly imbalanced decision problem.
  \item Compare supervised learning, anomaly detection, graph-based methods, and rule+ML hybrid approaches.
  \item Design a trust-and-safety workflow that combines scoring, thresholds, review operations, and business guardrails.
\end{itemize}

\section{Why fraud and trust \& safety matter}
E-commerce systems are open economic systems.
They make it easy for legitimate customers to browse, buy, receive promotions, request returns, and manage accounts.
Unfortunately, the same accessibility also creates opportunities for malicious behavior.

Fraud and abuse affect more than chargeback loss.
They influence:
\begin{itemize}
  \item payment costs,
  \item customer trust,
  \item seller trust in marketplaces,
  \item operational review burden,
  \item promotional efficiency,
  \item regulatory and compliance exposure.
\end{itemize}

For this reason, fraud prevention is best understood as part of a broader \emph{trust \& safety} capability rather than a narrow payment model.
The system must protect revenue while preserving a smooth user experience for legitimate customers.

\section{Fraud and abuse categories}
Fraud in e-commerce is heterogeneous.
Different abuse patterns require different signals and response strategies.

\subsection{Payment fraud}
Payment fraud often involves stolen cards, synthetic identities, mule accounts, or coordinated abuse of payment flows.
The platform must distinguish legitimate but unusual behavior from malicious transactions.

\subsection{Account takeover}
Account takeover occurs when an attacker gains control of an existing customer account.
This can lead to unauthorized purchases, loyalty theft, stored-value abuse, and privacy harm.

\subsection{Promotion and coupon abuse}
Promotions are common attack surfaces because they create direct economic value.
Abuse patterns include:
\begin{itemize}
  \item creating multiple accounts to redeem new-user offers,
  \item coordinated coupon sharing,
  \item automated checkout attempts,
  \item exploiting referral programs or cashback loops.
\end{itemize}

\subsection{Refund and return abuse}
Refund abuse includes false non-delivery claims, empty-box returns, repeated wardrobing behavior, and manipulation of return policies.
This category often requires joining transaction, logistics, support, and account-history signals.

\subsection{Marketplace and seller abuse}
In multi-seller environments, trust \& safety also includes:
\begin{itemize}
  \item counterfeit listings,
  \item fake reviews,
  \item seller collusion,
  \item policy evasion,
  \item manipulated performance metrics.
\end{itemize}

\section{Fraud as a machine learning problem}
Fraud detection is not a standard balanced classification task.
It has several defining properties.

\subsection{Class imbalance}
Fraudulent events are typically a small fraction of all transactions.
This means a naive high-accuracy classifier may still be useless if it mostly predicts ``not fraud.''

\subsection{Adversarial behavior}
Fraudsters adapt.
Once a platform changes a rule or scoring model, attackers may alter devices, payment instruments, shipping behavior, or timing patterns.
This makes fraud a moving-target problem.

\subsection{Label delay and label noise}
Ground truth often arrives late.
Chargebacks, disputes, and investigations may occur days or weeks after the original transaction.
Even then, labels can be incomplete or noisy.

\subsection{Asymmetric costs}
False negatives lose money and trust.
False positives block legitimate customers and harm conversion.
The decision threshold must therefore reflect business costs rather than abstract accuracy alone.

\section{Signals and features for fraud detection}
Fraud models typically depend on a broad feature space spanning account, transaction, device, and network behavior.

\subsection{Transaction signals}
\begin{itemize}
  \item order value,
  \item payment method,
  \item basket composition,
  \item discount depth,
  \item shipping speed,
  \item billing-shipping mismatch,
  \item repeated authorization attempts.
\end{itemize}

\subsection{Account signals}
\begin{itemize}
  \item account age,
  \item prior successful orders,
  \item return history,
  \item password reset activity,
  \item loyalty balance behavior,
  \item contact-information change frequency.
\end{itemize}

\subsection{Device and network signals}
\begin{itemize}
  \item device fingerprint,
  \item IP reputation,
  \item geolocation inconsistency,
  \item emulator or automation indicators,
  \item velocity across devices or sessions.
\end{itemize}

\subsection{Behavioral and sequence signals}
\begin{itemize}
  \item login and checkout timing,
  \item session depth,
  \item rapid account creation followed by high-value purchase,
  \item repeated failed payments,
  \item navigation patterns inconsistent with normal shopping behavior.
\end{itemize}

These signals are only useful when identity linkage and event quality are sound.
Fraud modeling therefore depends directly on the data-foundation chapter and on integration with operational systems.

\section{Rule-based systems}
Fraud prevention historically relied heavily on rules.
Rules remain important because they are fast, transparent, and easy to deploy for known abuse patterns.

\subsection{Examples of useful rules}
\begin{itemize}
  \item block transactions from known bad cards or IPs,
  \item trigger review if order value exceeds a threshold and account age is very low,
  \item limit coupon redemption frequency by account, device, and payment instrument,
  \item hold high-risk shipping destinations when signals conflict.
\end{itemize}

\subsection{Strengths and weaknesses}
Rules are valuable for:
\begin{itemize}
  \item deterministic policy enforcement,
  \item emergency response,
  \item explainability in operations.
\end{itemize}

They are weak when:
\begin{itemize}
  \item abuse patterns evolve quickly,
  \item interactions among many weak signals matter,
  \item manual rule maintenance becomes operationally expensive.
\end{itemize}

\section{Supervised learning}
Supervised models learn from labeled examples of fraud and legitimate behavior.
They are widely used because they can combine many signals more flexibly than hand-written rules.

\subsection{Typical model families}
Common model families include:
\begin{itemize}
  \item logistic regression,
  \item tree-based ensembles,
  \item gradient boosting,
  \item neural models where scale and feature complexity justify them.
\end{itemize}

\subsection{Why supervised learning helps}
Supervised learning helps by:
\begin{itemize}
  \item aggregating many weak predictors,
  \item ranking risk more smoothly than rigid rules,
  \item enabling threshold tuning by cost and segment,
  \item adapting to new feature sets over time.
\end{itemize}

\subsection{Operational caution}
Supervised models are only as good as:
\begin{itemize}
  \item the quality of labels,
  \item the recency of training data,
  \item the alignment between training and serving features,
  \item the monitoring for drift and policy changes.
\end{itemize}

\section{Anomaly detection}
Some abuse patterns are rare or novel enough that labeled examples are limited.
In these settings, anomaly detection can be useful.

\subsection{What anomaly methods look for}
Anomaly-oriented methods try to identify behavior that differs sharply from established norms, such as:
\begin{itemize}
  \item unusual transaction timing,
  \item atypical device usage,
  \item unexpected country or route changes,
  \item extreme coupon redemption behavior,
  \item sudden seller or account behavior shifts.
\end{itemize}

\subsection{Strengths and weaknesses}
Anomaly methods are useful for surfacing suspicious patterns and complementing supervised systems.
However, anomaly does not automatically mean fraud.
Legitimate edge cases can also appear unusual.
This means anomaly scores usually work best as part of a broader review or hybrid decision pipeline.

\section{Graph-based fraud analysis}
Fraud is often relational.
Accounts, devices, cards, addresses, merchants, and products form networks rather than isolated rows of data.

\subsection{Why graph structure matters}
Graph methods help detect:
\begin{itemize}
  \item many accounts linked to one device,
  \item many cards shipping to related addresses,
  \item collusive seller or review networks,
  \item abuse rings coordinating promotions or refunds.
\end{itemize}

\subsection{Graph features and graph ML}
Practical graph-based approaches may include:
\begin{itemize}
  \item connected-component and degree features,
  \item shared-entity counts,
  \item neighborhood risk propagation,
  \item graph neural networks where scale and maturity justify them.
\end{itemize}

Graph analysis is powerful because fraudsters often try to hide in transactional noise while still reusing infrastructure or identity fragments.

\section{Hybrid rules + ML systems}
In production e-commerce systems, the strongest approach is often hybrid rather than purely rule-based or purely model-based.

\subsection{Why hybrids work}
Rules are good for hard policy constraints and urgent pattern blocking.
ML is good for scoring complex combinations of soft signals.
Combined systems can:
\begin{itemize}
  \item enforce mandatory blocks,
  \item assign risk scores,
  \item route borderline cases to review,
  \item and learn from review outcomes.
\end{itemize}

\subsection{Example decision flow}
A practical decision flow may look like:
\begin{enumerate}
  \item apply deterministic policy rules,
  \item compute risk score from supervised and anomaly signals,
  \item adjust using graph or network intelligence,
  \item compare against thresholds for approve, review, or reject,
  \item log explanations and outcomes for feedback.
\end{enumerate}

\section{Thresholds and decision economics}
Fraud systems do not end at model scoring.
The business must choose action thresholds.

\subsection{Approve, review, reject}
Many systems use three zones:
\begin{itemize}
  \item \textbf{Approve:} low-risk transactions proceed automatically.
  \item \textbf{Review:} medium-risk cases go to human review or secondary checks.
  \item \textbf{Reject or challenge:} high-risk cases are blocked or subjected to stronger verification.
\end{itemize}

\subsection{Cost-sensitive tuning}
Thresholds should reflect:
\begin{itemize}
  \item fraud loss,
  \item customer lifetime value,
  \item review-team capacity,
  \item conversion impact,
  \item chargeback and compliance costs.
\end{itemize}

This means threshold selection is not a pure ML exercise.
It is an economic operating decision.

\section{Human-in-the-loop review}
Not all fraud decisions should be fully automated.
Human review remains important for edge cases, novel fraud patterns, appeals, and policy-sensitive decisions.

\subsection{What reviewers need}
Review tools should expose:
\begin{itemize}
  \item transaction history,
  \item device and network context,
  \item linked entities,
  \item model score and contributing signals,
  \item policy reason codes,
  \item prior outcomes on related entities.
\end{itemize}

\subsection{Feedback loop}
Review outcomes should feed back into:
\begin{itemize}
  \item rule refinement,
  \item model retraining,
  \item investigation workflows,
  \item reporting on false positives and operational load.
\end{itemize}

\section{Fraud evaluation metrics}
Accuracy alone is misleading in fraud detection because the class distribution is highly imbalanced.

\subsection{Useful metrics}
Practical evaluation often includes:
\begin{itemize}
  \item precision,
  \item recall,
  \item precision-recall curves,
  \item ROC-AUC with caution,
  \item false-positive rate,
  \item fraud-dollar capture,
  \item review rate,
  \item cost-weighted utility.
\end{itemize}

\subsection{Business interpretation}
A model with strong recall but excessive false positives may block good customers.
A model with low review cost but weak fraud capture may lose large sums.
Evaluation must therefore align with operational and economic goals.

\section{Drift, adaptation, and monitoring}
Fraud systems need continuous monitoring because attacker behavior evolves.

\subsection{What to monitor}
\begin{itemize}
  \item score distribution shifts,
  \item feature drift,
  \item approval and decline rates,
  \item review queue volume,
  \item chargeback lag trends,
  \item abuse-type mix over time.
\end{itemize}

\subsection{Why monitoring is essential}
Even a well-trained model can degrade if:
\begin{itemize}
  \item fraud patterns shift,
  \item product flows change,
  \item new promotions alter normal behavior,
  \item identity or logging systems are modified.
\end{itemize}

Fraud detection is therefore inseparable from MLOps and operational analytics.

\section{Trust \& safety beyond payment fraud}
Trust \& safety is broader than payment authorization.
It includes keeping the entire commerce environment healthy.

\subsection{Examples beyond transactions}
\begin{itemize}
  \item review authenticity,
  \item marketplace seller compliance,
  \item abuse of support or refund workflows,
  \item harmful or prohibited catalog content,
  \item bot-driven scraping or stock hoarding.
\end{itemize}

\subsection{Why this matters architecturally}
These problems span domains.
A trust \& safety capability often depends on:
\begin{itemize}
  \item commerce events,
  \item ERP refund records,
  \item CRM support interactions,
  \item identity and authentication systems,
  \item marketplace governance workflows.
\end{itemize}

This reinforces the enterprise view that trust \& safety is an operating system for commerce, not a single checkout model.

\section{Hands-on lab: risk scoring and thresholding}
The practical goal of this chapter is to connect fraud modeling with decision policy.

\subsection{Lab scenario}
Students are given a synthetic transaction dataset containing payment, account, device, and label information, plus estimated fraud-loss and false-positive costs.

\subsection{Lab tasks}
\begin{enumerate}
  \item Build a baseline risk model using structured transaction and account features.
  \item Compute precision, recall, and a cost-sensitive evaluation summary.
  \item Define approve, review, and reject thresholds.
  \item Propose one anomaly signal and one graph-style linkage feature to improve the system.
  \item Explain how a human review queue would be integrated into the workflow.
\end{enumerate}

\subsection{Expected learning outcome}
By the end of the lab, students should be able to explain fraud detection as a staged decision system rather than a single prediction score.

\section{Managerial implications}
Fraud and trust \& safety systems require balancing protection with customer experience.
Overly strict controls reduce conversion and frustrate legitimate users.
Overly weak controls invite direct loss and operational abuse.

The best systems therefore combine:
\begin{itemize}
  \item clear policy,
  \item strong data quality,
  \item hybrid detection methods,
  \item review operations,
  \item and continuous monitoring.
\end{itemize}

Organizations that treat fraud as an isolated technical model usually underinvest in the workflow and governance layers that actually determine business outcomes.

\section{Multiple-choice questions (MCQs)}
\begin{enumerate}
  \item Why is fraud detection in e-commerce often difficult from an ML perspective?
  \begin{enumerate}
    \item Because fraud labels are always complete and immediate.
    \item Because the problem is highly imbalanced, adversarial, and often affected by delayed labels.
    \item Because fraud can only be solved with static rules.
    \item Because customer behavior never changes.
  \end{enumerate}

  \item Which is the best example of \emph{promotion abuse}?
  \begin{enumerate}
    \item A customer browsing many products before purchase.
    \item Creating many accounts or identities to repeatedly exploit a first-order coupon.
    \item A warehouse updating shipment status.
    \item A user selecting a shipping address.
  \end{enumerate}

  \item Why are graph-based methods useful in fraud detection?
  \begin{enumerate}
    \item Because fraud is often relational and linked through shared devices, addresses, accounts, or payment instruments.
    \item Because graphs eliminate the need for labels.
    \item Because they only work for product recommendation.
    \item Because they remove the need for monitoring.
  \end{enumerate}

  \item In a hybrid fraud system, rules are especially useful for:
  \begin{enumerate}
    \item hard policy enforcement and fast blocking of known bad patterns.
    \item replacing all human review.
    \item guaranteeing perfect recall.
    \item making embeddings unnecessary in every domain.
  \end{enumerate}

  \item Why are thresholds in fraud systems business decisions as well as technical decisions?
  \begin{enumerate}
    \item Because thresholds determine the tradeoff between fraud loss, false positives, review cost, and customer experience.
    \item Because thresholds only affect documentation style.
    \item Because all thresholds are set randomly.
    \item Because supervised learning does not produce scores.
  \end{enumerate}
\end{enumerate}

\subsection*{Answer key}
\begin{enumerate}
  \item (b)
  \item (b)
  \item (a)
  \item (a)
  \item (a)
\end{enumerate}

\section{Exercises (short)}
\begin{enumerate}
  \item Compare rule-based detection, supervised learning, and anomaly detection for one fraud scenario of your choice. State one advantage and one limitation of each.
  \item Design a minimal feature set for detecting account takeover. Include account, device, sequence, and review-history signals.
  \item Define a three-zone approve/review/reject threshold policy for a checkout fraud model and explain how review-team capacity changes your choice.
\end{enumerate}

\section{Mini-case (Odoo-linked, vendor-neutral)}
\textbf{Scenario:} A hybrid commerce company uses a headless storefront, Odoo for orders and refunds, and a central data platform. The business is seeing rising chargebacks, coupon abuse, and suspicious refund patterns.

\textbf{Task:} Design a trust \& safety workflow:
\begin{itemize}
  \item Define which events and records must be linked across storefront, payment, Odoo, and support systems.
  \item Propose a hybrid detection design using rules, ML scoring, and at least one graph-style signal.
  \item Specify the criteria for automatic approval, manual review, and automatic rejection in one high-risk flow such as refund approval or checkout authorization.
\end{itemize}
